% Part: first-order-logic
% Chapter: tableaux
% Section: soundness.tex

\documentclass[../../../include/open-logic-section]{subfiles}

\begin{document}

\iftag{FOL}
      {\olfileid{fol}{tab}{sou}}
      {\olfileid{pl}{tab}{sou}}
      
\olsection{Correctheid: het systeem leidt niets onjuists af}

\begin{explain}
!!^a{derivation}ssysteem, zoals de tableaumethode, is \emph{correct}
als het niets kan !!{derive} wat niet echt geldt. Correctheid geeft dus
een veiligheidsgarantie voor !!{derivation}ssystemen. Welke garantie
we precies nodig hebben, hangt af van de bewijstheoretische eigenschap
waarover het gaat. We willen bijvoorbeeld weten dat
\begin{enumerate}
\item elke !!{derivable}~$!A$ ook echt
  \iftag{FOL}{geldig}{een tautologie} is;
\item als !!a{sentence} formeel uit andere zinnen !!{derivable} is, hij
  ook logisch uit die zinnen volgt;
\item als een verzameling !!{sentence}s inconsistent is, zij
  onvervulbaar is: de zinnen zijn niet allemaal tegelijk waar te maken.
\end{enumerate}
Deze eigenschappen zijn belangrijk. Als een ervan niet geldt, deugt
het !!{derivation}ssysteem niet: het zou te veel !!{derive}. Daarom
moeten we bewijzen dat zo'n systeem correct is.

Alle drie eigenschappen hierboven zijn met gesloten !!{tableau}s
gedefinieerd. Om (1)--(3) te bewijzen, moeten we daarom iets laten zien
over de betekenis van zo'n gesloten !!{tableau}. Eerst leggen we vast
wanneer een structuur !!a{signed formula} waar maakt. Daarna bewijzen
we dat geen structuur alle aannamen van een gesloten !!{tableau}
tegelijk waar maakt. Daaruit volgen (1)--(3) meteen als gevolgen.
\end{explain}

\begin{defn}
\iftag{FOL}{!!^a{structure}}{!!^a{valuation}}~$\iftag{FOL}{\Struct{M}}{\pAssign{v}}$
\emph{maakt} !!a{signed formula} $\sFmla{\True}{!A}$ waar
precies als $\iftag{FOL}{\Sat{M}}{\pSat{v}}{!A}$. Deze maakt
$\sFmla{\False}{!A}$ waar precies als
$\iftag{FOL}{\Sat/{M}}{\pSat/{v}}{!A}$. $\iftag{FOL}{\Struct{M}}{\pAssign{v}}$
maakt een verzameling !!{signed formula}s~$\Gamma$ waar precies als
elke $\sFmla{S}{!A} \in \Gamma$ waar wordt gemaakt. $\Gamma$~is
\emph{vervulbaar} als er een
\iftag{FOL}{!!a{structure}}{!!a{valuation}} bestaat die alle formules
erin waar maakt. Anders is de verzameling \emph{onvervulbaar}.
\end{defn}

\begin{thm}[Correctheid]
  \ollabel{thm:tableau-soundness}
  Als $\Gamma$ een gesloten !!{tableau} heeft, is $\Gamma$
  onvervulbaar.
\end{thm}

\begin{proof}
Noem een tak van !!a{tableau} vervulbaar als alle !!{signed formula}s
op die tak samen vervulbaar zijn. Noem !!a{tableau} vervulbaar als het
minstens één vervulbare tak heeft.

We bewijzen eerst één bewering. Breid een vervulbaar !!{tableau} met
een afleidingsregel uit. Dan is ook het nieuwe !!{tableau} vervulbaar.
Daaruit volgt de stelling. Elk gesloten !!{tableau} ontstaat namelijk
door afleidingsregels toe te passen op het !!{tableau} dat alleen
aannamen uit~$\Gamma$ bevat. Als $\Gamma$ vervulbaar was, zou elk
!!{tableau} ervoor dus vervulbaar blijven. Maar een gesloten
!!{tableau} is niet vervulbaar. Elke tak bevat zowel
$\sFmla{\True}{!A}$ als $\sFmla{\False}{!A}$. Geen structuur kan~$!A$
tegelijk wel en niet waar maken.

Neem een vervulbaar !!{tableau}. Het heeft dus minstens één vervulbare
tak. Een afleidingsregel doet een van twee dingen: hij voegt
!!{signed formula}s aan een tak toe, of hij splitst een tak in tweeën.
Misschien breidt de regel een andere tak uit dan een vervulbare tak.
Dan blijft die vervulbare tak gewoon bestaan in het nieuwe
!!{tableau}. Ook het nieuwe tableau is dan vervulbaar. We hoeven dus
alleen te kijken naar het geval waarin de regel juist op een
vervulbare tak wordt toegepast.

Laat $\Gamma$ de verzameling !!{signed formula}s op die tak zijn. Laat
$\sFmla{S}{!A} \in \Gamma$ de !!{signed formula} zijn waarop we de
regel toepassen. Splitst de regel de tak niet? Dan moeten we laten zien
dat de uitgebreide tak nog steeds vervulbaar is. Die tak bestaat uit
$\Gamma$ samen met de conclusies van de regel. Splitst de regel de tak
wel? Dan moeten we laten zien dat minstens één van de twee nieuwe
takken vervulbaar is.
  
Eerst bekijken we de regels die de tak niet splitsen.
\begin{enumerate}
\item Pas $\TRule{\True}{\lnot}$ toe op
  $\sFmla{\True}{\lnot !B} \in \Gamma$. De uitgebreide tak bevat dan
  de !!{signed formula}s $\Gamma \cup
  \{\sFmla{\False}{!B}\}$. Stel dat
  $\iftag{FOL}{\Sat{M}}{\pSat{v}}{\Gamma}$. Dan geldt in het bijzonder
  $\iftag{FOL}{\Sat{M}}{\pSat{v}}{\lnot !B}$. Dus geldt
  $\iftag{FOL}{\Sat/{M}}{\pSat/{v}}{!B}$. Dat wil zeggen:
  $\iftag{FOL}{\Struct{M}}{\pAssign{v}}$ maakt
  $\sFmla{\False}{!B}$.
\item Pas $\TRule{\False}{\lnot}$ toe op
  $\sFmla{\False}{\lnot !B} \in \Gamma$: opgave.
\item Pas $\TRule{\True}{\land}$ toe op
  $\sFmla{\True}{!B \land !C} \in \Gamma$. Daardoor komen er twee
  nieuwe !!{signed formula}s op de tak: $\sFmla{\True}{!B}$ en
  $\sFmla{\True}{!C}$. Stel dat
  $\iftag{FOL}{\Sat{M}}{\pSat{v}}{\Gamma}$. Dan geldt in het bijzonder
  $\iftag{FOL}{\Sat{M}}{\pSat{v}}{!B \land !C}$. Daaruit volgen
  $\iftag{FOL}{\Sat{M}}{\pSat{v}}{!B}$ en
  $\iftag{FOL}{\Sat{M}}{\pSat{v}}{!C}$. Dit betekent dat
  $\iftag{FOL}{\Struct{M}}{\pAssign{v}}$ zowel
  $\sFmla{\True}{!B}$ als $\sFmla{\True}{!C}$ waar maakt.
\item Pas $\TRule{\False}{\lor}$ toe op
  $\sFmla{\False}{!B \lor !C} \in \Gamma$: opgave.
\item Pas $\TRule{\False}{\lif}$ toe op
  $\sFmla{\False}{!B \lif !C} \in \Gamma$. Daardoor komen er twee
  nieuwe !!{signed formula}s op de tak: $\sFmla{\True}{!B}$ en
  $\sFmla{\False}{!C}$. Stel dat
  $\iftag{FOL}{\Sat{M}}{\pSat{v}}{\Gamma}$. Dan geldt in het bijzonder
  $\iftag{FOL}{\Sat/{M}}{\pSat/{v}}{!B \lif !C}$. Daaruit volgen
  $\iftag{FOL}{\Sat{M}}{\pSat{v}}{!B}$ en
  $\iftag{FOL}{\Sat/{M}}{\pSat/{v}}{!C}$. Dit betekent dat
  $\iftag{FOL}{\Struct{M}}{\pAssign{v}}$ zowel
  $\sFmla{\True}{!B}$ als $\sFmla{\False}{!C}$ waar maakt.

\iftag{FOL}{%
\item Pas $\TRule{\True}{\lforall}$ toe op
  $\sFmla{\True}{\lforall[x][!B(x)]} \in \Gamma$. Daardoor komt de
  nieuwe !!{signed formula}~$\sFmla{\True}{!A(t)}$ op de tak. Stel dat
  $\Sat{M}{\Gamma}$. Dan geldt in het bijzonder
  $\Sat{M}{\lforall[x][!A(x)]}$. Volgens
  \olref[syn][sem]{prop:quant-terms} geldt dan $\Sat{M}{!A(t)}$. Dus
  maakt $\Struct{M}$ de formule $\sFmla{\True}{!A(t)}$ waar.

\item Pas $\TRule{\False}{\lforall}$ toe op
  $\sFmla{\False}{\lforall[x][!B(x)]} \in \Gamma$. Daardoor komt de
  nieuwe !!{signed formula}~$\sFmla{\False}{!A(a)}$ op de tak. Hier is
  $a$ !!a{constant} die niet in~$\Gamma$ voorkomt. Omdat $\Gamma$
  vervulbaar is, bestaat er een $\Struct{M}$ waarvoor
  $\Sat{M}{\Gamma}$. In het bijzonder geldt dan
  $\Sat/{M}{\lforall[x][!B(x)]}$. We moeten laten zien dat
  $\Gamma \cup \{\sFmla{\False}{!A(a)}\}$ vervulbaar is. Daarvoor
  maken we als volgt een geschikte~$\Struct{M'}$.

  Volgens \olref[syn][ass]{prop:sat-quant} geldt
  $\Sat/{M}{\lforall[x][!B(x)]}$ precies als er een $s$ is waarvoor
  $\Sat/{M}{!B(x)}[s]$. Neem nu een $\Struct{M'}$ die hetzelfde is als
  $\Struct{M}$, behalve dat $\Assign{a}{M'} = s(x)$. Volgens
  \olref[syn][ext]{cor:extensionality-sent} geldt voor elke
  $\sFmla{\True}{!C} \in \Gamma$ dat $\Sat{M'}{!C}$. Voor elke
  $\sFmla{\False}{!C} \in \Gamma$ geldt dat $\Sat/{M'}{!C}$. Dat komt
  doordat $a$ niet in~$\Gamma$ voorkomt.

  Volgens \olref[syn][ext]{prop:extensionality} geldt
  $\Sat/{M'}{!A(x)}[s]$. Volgens
  \olref[syn][ext]{prop:ext-formulas} geldt vervolgens
  $\Sat/{M'}{!A(a)}[s]$. Omdat $!A(a)$ !!a{sentence} is, volgt uit
  \olref[syn][ass]{prop:sentence-sat-true} dat
  $\Sat/{M'}{!A(a)}$. Dat wil zeggen: $\Struct{M'}$ maakt
  $\sFmla{\False}{!A(a)}$ waar.
\item Pas $\TRule{\True}{\lexists}$ toe op
  $\sFmla{\True}{\lexists[x][!B(x)]} \in \Gamma$: opgave.
\item Pas $\TRule{\False}{\lexists}$ toe op
  $\sFmla{\False}{\lexists[x][!B(x)]} \in \Gamma$: opgave.}{}
\end{enumerate}
Nu bekijken we de regels die een tak wel splitsen.
\begin{enumerate}
\item Pas $\TRule{\False}{\land}$ toe op
  $\sFmla{\False}{!B \land !C} \in \Gamma$. Daardoor ontstaan twee
  takken. De linkertak gaat verder met $\sFmla{\False}{!B}$ en de
  rechtertak met $\sFmla{\False}{!C}$. Stel dat
  $\iftag{FOL}{\Sat{M}}{\pSat{v}}{\Gamma}$. Dan geldt in het bijzonder
  $\iftag{FOL}{\Sat/{M}}{\pSat/{v}}{!B \land !C}$. Daaruit volgt
  $\iftag{FOL}{\Sat/{M}}{\pSat/{v}}{!B}$ of
  $\iftag{FOL}{\Sat/{M}}{\pSat/{v}}{!C}$. In het eerste geval maakt
  $\iftag{FOL}{\Struct{M}}{\pAssign{v}}$
  $\sFmla{\False}{!B}$ waar. Daarmee maakt
  $\iftag{FOL}{\Struct{M}}{\pAssign{v}}$ alle formules op de linkertak
  waar. In het tweede geval maakt
  $\iftag{FOL}{\Struct{M}}{\pAssign{v}}$
  $\sFmla{\False}{!C}$ waar. Daarmee maakt
  $\iftag{FOL}{\Struct{M}}{\pAssign{v}}$ alle formules op de
  rechtertak waar.
\item Pas $\TRule{\True}{\lor}$ toe op
  $\sFmla{\True}{!B \lor !C} \in \Gamma$: opgave.
\item Pas $\TRule{\True}{\lif}$ toe op
  $\sFmla{\True}{!B \lif !C} \in \Gamma$: opgave.
\item Breid de tak uit met \Cut. Daardoor ontstaan twee takken. De ene
  bevat $\sFmla{\True}{!B}$ en de andere
  $\sFmla{\False}{!B}$. We weten dat
  $\iftag{FOL}{\Sat{M}}{\pSat{v}}{\Gamma}$ en dat ofwel
  $\iftag{FOL}{\Sat{M}}{\pSat{v}}{!B}$, ofwel
  $\iftag{FOL}{\Sat/{M}}{\pSat/{v}}{!B}$. Daarom maakt
  $\iftag{FOL}{\Struct{M}}{\pAssign{v}}$ de linker- of de rechtertak
  waar.
\end{enumerate}
\end{proof}

\tagprob{FOL}
\begin{prob}
Maak het bewijs van \olref[fol][tab][sou]{thm:tableau-soundness} af.
\end{prob}
\tagendprob

\tagprob{notFOL}
\begin{prob}
Maak het bewijs van \olref[pl][tab][sou]{thm:tableau-soundness} af.
\end{prob}
\tagendprob

\begin{cor}
\ollabel{cor:weak-soundness}
Als $\Proves !A$, dan is $!A$ \iftag{FOL}{geldig}{een tautologie}.
\end{cor}

\begin{cor}
\ollabel{cor:entailment-soundness}
Als $\Gamma \Proves !A$, dan geldt ook $\Gamma \Entails !A$.
\end{cor}

\begin{proof}
  Stel dat $\Gamma \Proves !A$. Dan zijn er $!B_1$, \dots,
  $!B_n \in \Gamma$ waarvoor
  $\{\sFmla{\False}{!A}, \sFmla{\True}{!B_1}, \dots,
  \sFmla{\True}{!B_n}\}$ een gesloten !!{tableau} heeft. Volgens
  \olref{thm:tableau-soundness} maakt elke
  \iftag{FOL}{!!{structure}}{!!{valuation}}~$\iftag{FOL}{\Struct{M}}{\pAssign{v}}$
  een zekere $!B_i$ onwaar of $!A$ waar. Als dus
  $\iftag{FOL}{\Sat{M}}{\pSat{v}}{\Gamma}$, dan geldt ook
  $\iftag{FOL}{\Sat{M}}{\pSat{v}}{!A}$.
\end{proof}

\begin{cor}
\ollabel{cor:consistency-soundness}
Als $\Gamma$ vervulbaar is, dan is zij consistent.
\end{cor}

\begin{proof}
We bewijzen de contrapositie. Neem aan dat $\Gamma$ niet consistent
is. Dan zijn er $!B_1$, \dots, $!B_n \in \Gamma$ en is er een gesloten
!!{tableau} voor
$\{\sFmla{\True}{!B_1}, \dots, \sFmla{\True}{!B_n}\}$. Volgens
\olref{thm:tableau-soundness} bestaat er geen
$\iftag{FOL}{\Struct{M}}{\pAssign{v}}$ waarvoor
$\iftag{FOL}{\Sat{M}}{\pSat{v}}{!B_i}$ voor alle $i=1$, \dots,~$n$.
Daarom is $\Gamma$ niet vervulbaar.
\end{proof}

\end{document}
